% ID: MAT-TRI-EQG-006
% Titolo: Intersezioni di una funzione goniometrica con gli assi (3)
% Disciplina: Matematica
% Area: trigonometria
% Argomento: equazioni_goniometriche
% Classe: Quarta secondaria di secondo grado
% Difficolta: 4
% Tipo: grafico
% Tag: matematica, trigonometria, equazioni_goniometriche
% Fonte: verifica_goniometria_avanzata_anonimizzata
% Provenienza: materiale fornito dal docente; testo anonimizzato; soluzione redatta per Exergo

\begin{esercizio}
Trova le coordinate di tutte le intersezioni con gli assi
cartesiani del grafico della funzione: 
\[
f(x)=\cos 3x+\cos x-\sqrt{3}\cos 2x.
\]
Considera $x\in[0;2\pi)$.
Salvo diversa indicazione, lavora in $\mathbb R$; angoli in radianti.
\end{esercizio}

\begin{soluzione}
Usando solo duplicazione e addizione,
$\cos3x=\cos(2x+x)=(2\cos^2x-1)\cos x-2\sin^2x\cos x=4\cos^3x-3\cos x$.
Quindi $f(x)=(2\cos x-\sqrt3)(2\cos^2x-1)$.
Le intersezioni con l'asse $x$ sono $(a,0)$ per
$a\in\{\pi/6,\pi/4,3\pi/4,5\pi/4,7\pi/4,11\pi/6\}$.
Con l'asse $y$: $(0,2-\sqrt3)$.
\end{soluzione}
